% !TeX root = ../main.tex
\chapter{HƯỚNG DẪN TÁI TẠO KẾT QUẢ}
\label{app:D}

\section{Cấu trúc kho mã nguồn}

\begin{table}[H]
\centering
\caption{Cấu trúc thư mục chính của kho mã nguồn}
\label{tab:app_repo}
\begin{tabular}{L{5.2cm} L{7.8cm}}
\toprule
\textbf{Đường dẫn} & \textbf{Nội dung} \\
\midrule
\texttt{notebooks/} & Mười \en{notebook} thực nghiệm, đánh số theo thứ tự chạy \\
\texttt{notebooks/\allowbreak data/\allowbreak raw/} & Ảnh gốc, chia sẵn theo tập và theo lớp \\
\texttt{notebooks/\allowbreak data/\allowbreak pseudo\_\allowbreak labels*/} & Các bộ nhãn giả sinh ra \\
\texttt{notebooks/\allowbreak models/} & Điểm kiểm tra và tệp chỉ số của các mô hình \\
\texttt{utils/} & Mã dùng chung: mô hình, tiền xử lý, chỉ số, Grad-CAM \\
\texttt{visualizations/} & Hình sinh ra từ các \en{notebook} \\
\texttt{app.py} & Ứng dụng minh hoạ Streamlit \\
\texttt{docs/thesis/} & Mã nguồn LaTeX của luận văn \\
\bottomrule
\end{tabular}
\end{table}

\section{Thứ tự thực thi}

Các \en{notebook} phải chạy theo đúng thứ tự vì mỗi bước sử dụng sản phẩm của bước
trước.

\begin{table}[H]
\centering
\caption{Thứ tự thực thi và sản phẩm của từng notebook}
\label{tab:app_order}
\small
\begin{tabular}{c L{4.6cm} L{6.4cm}}
\toprule
\textbf{STT} & \textbf{Notebook} & \textbf{Sản phẩm} \\
\midrule
1 & \texttt{01-\allowbreak data-\allowbreak exploration} & Hình EDA, thống kê bộ dữ liệu \\
2 & \texttt{02-\allowbreak classification-\allowbreak baseline} & \texttt{best\_\allowbreak model.\allowbreak pth}, chỉ số phân loại \\
3 & \texttt{03-\allowbreak xai-\allowbreak gradcam-\allowbreak analysis} & Hình bản đồ kích hoạt, thống kê độ tin cậy \\
4 & \texttt{04-\allowbreak pseudo-\allowbreak labeling} & 3.104 mặt nạ nhãn giả V1 \\
5 & \texttt{05-\allowbreak segmentation-\allowbreak model} & Mô hình phân đoạn V1, chỉ số kiểm thử \\
6 & \texttt{06-\allowbreak pseudo-\allowbreak labeling-\allowbreak v2} & 3.104 mặt nạ nhãn giả V2 \\
7 & \texttt{07-\allowbreak segmentation-\allowbreak v2} & Mô hình phân đoạn V2, chỉ số kiểm thử \\
8 & \texttt{08-\allowbreak segmentation-\allowbreak v3-\allowbreak sam} & Nhãn tinh chỉnh SAM, mô hình V3 \\
9 & \texttt{10-\allowbreak inference-\allowbreak single-\allowbreak image} & Hình minh hoạ suy luận đầu cuối \\
\bottomrule
\end{tabular}
\end{table}

Riêng \en{notebook} thứ tám có cơ chế tiếp tục: nếu thư mục nhãn tinh chỉnh đã đủ
3.104 tệp, bước sinh nhãn bằng \SAM{} sẽ được bỏ qua và mô hình \SAM{} không được
nạp lên bộ nhớ đồ hoạ. Muốn sinh lại từ đầu cần đặt biến \texttt{FORCE\_REDO} bằng
\texttt{True} hoặc xoá thư mục nhãn.

Các \en{notebook} huấn luyện đều lưu điểm kiểm tra tạm sau mỗi epoch, cho phép
tiếp tục nếu quá trình bị gián đoạn.

\section{Yêu cầu môi trường}

Môi trường đã dùng để sinh toàn bộ kết quả trong luận văn được liệt kê ở
\cref{tab:moitruong}. Các thư viện chính cần cài đặt:

\begin{itemize}
    \item \texttt{torch}, \texttt{torchvision} --- phiên bản tương thích CUDA của
          máy đích.
    \item \texttt{segmentation-\allowbreak models-\allowbreak pytorch} --- cung cấp kiến trúc \UNetpp{}.
    \item \texttt{albumentations} --- tăng cường dữ liệu đồng bộ ảnh và mặt nạ.
    \item \texttt{segment-\allowbreak anything} --- mô hình \SAM{}.
    \item \texttt{opencv-\allowbreak python}, \texttt{pillow}, \texttt{numpy},
          \texttt{pandas}, \texttt{matplotlib}, \texttt{seaborn},
          \texttt{scikit-learn}, \texttt{tqdm}.
    \item \texttt{streamlit}, \texttt{plotly} --- chỉ cần cho ứng dụng minh hoạ.
\end{itemize}

Trọng số \SAM{} ViT-B (khoảng 375 MB) được tải tự động ở lần chạy đầu tiên của
\en{notebook} thứ tám.

\section{Lưu ý về khả năng tái lập}

\subsection{Các yếu tố đã được cố định}

Hạt giống ngẫu nhiên được đặt bằng 42 cho việc chia tập dữ liệu ở mọi phiên bản
phân đoạn, bảo đảm các phiên bản V2 và V3 dùng đúng cùng một tập kiểm thử.

\subsection{Các yếu tố chưa được cố định hoàn toàn}

Cần lưu ý ba nguồn sai khác có thể khiến kết quả chạy lại không trùng khớp tuyệt
đối với số liệu báo cáo:

\begin{itemize}
    \item \textbf{Tính không tất định của cuDNN.} Chỉ \en{notebook} thứ tám đặt
          \texttt{cudnn.\allowbreak deterministic} bằng \texttt{True}; các \en{notebook} còn
          lại không đặt. Các thuật toán tích chập được chọn tự động có thể khác
          nhau giữa các lần chạy và giữa các loại GPU.
    \item \textbf{Lấy mẫu điểm gợi ý cho \SAM{}.} Bốn điểm tiền cảnh và ba điểm
          nền được lấy ngẫu nhiên mà không cố định hạt giống, nên nhãn tinh chỉnh
          sinh lại sẽ khác đôi chút so với bộ nhãn đã dùng.
    \item \textbf{Tăng cường dữ liệu ngẫu nhiên.} Chuỗi phép biến đổi Albumentations
          không cố định hạt giống trong vòng lặp huấn luyện.
    \end{itemize}

Vì vậy, khi chạy lại, các chỉ số có thể lệch ở chữ số thập phân thứ hai hoặc thứ
ba. Các kết luận định tính của luận văn --- chiều của các so sánh, hiện tượng nở
mặt nạ, sự lệch pha giữa hàm mất mát và IoU --- không phụ thuộc vào mức sai khác
này.

\subsection{Khuyến nghị cho lần chạy tiếp theo}

Để tăng khả năng tái lập, ba thay đổi nên được áp dụng: cố định hạt giống ngẫu
nhiên cho \texttt{random}, \texttt{numpy} và \texttt{torch} ở đầu mọi
\en{notebook}; đặt \texttt{cudnn.\allowbreak deterministic} bằng \texttt{True} một cách nhất
quán; và cố định hạt giống cho bước lấy mẫu điểm gợi ý của \SAM{}.

\section{Truy vết số liệu}

Toàn bộ con số xuất hiện trong luận văn được ghi lại trong tệp
\texttt{docs/\allowbreak thesis/\allowbreak INVENTORY.\allowbreak md}, kèm thông tin về \en{notebook} và ô lệnh đã
sinh ra nó. Tệp này cũng liệt kê các điểm mà số liệu chưa đủ bằng chứng, cùng
những chỗ số liệu trong công bố liên quan của nhóm không khớp với kết quả thực
thi thực tế.

Một số giá trị trong luận văn được tính lại trực tiếp từ các tệp kết quả trên đĩa
thay vì lấy từ ô kết quả của \en{notebook} --- cụ thể là bảng độ phủ của nhãn tinh
chỉnh bằng \SAM{} ở \cref{tab:sam_expansion}. Đoạn mã dùng để tính lại được ghi
đầy đủ trong tệp kiểm kê nói trên, cho phép kiểm chứng độc lập.
